\section{Hierarchical Image Classification}

Hierarchical image classification \cite{silla2011survey,chang2021your,jiang2024hierarchical} extends standard image classification by predicting labels according to a predefined class hierarchy. Instead of assigning each image to only one flat class, a hierarchical classifier can produce predictions at multiple semantic levels, such as order, family, and species in bird recognition. This setting is useful when the label space naturally follows a taxonomy and when different levels of granularity are meaningful for different users or application scenarios.


In the context of this thesis, hierarchical classification is particularly relevant because bird categories are naturally organized taxonomically. A model that predicts only the species label may achieve fine-grained recognition, but it does not explicitly evaluate whether the predicted species is compatible with the corresponding family and order. Therefore, hierarchical classification provides a more structured formulation for bird recognition, where both fine-grained accuracy and taxonomy-level consistency are important.

\subsection{Flat Classification vs. Hierarchical Classification}

Flat classification treats all target categories as independent classes at a single semantic level \cite{silla2011survey}. In fine-grained bird classification, this usually means predicting one of the 200 species in CUB-200-2011. This formulation is simple and widely used, but it ignores the hierarchical relationships between classes. For example, if a model predicts a species from one family while its family-level prediction indicates a different family, the prediction path becomes inconsistent with the taxonomy.

Hierarchical classification, by contrast, explicitly models labels at multiple levels of abstraction. In a bird taxonomy, the model may predict the order, family, and species of an image. This formulation allows the model to represent both coarse and fine semantic information. It also makes evaluation more informative, since hierarchy-aware evaluation treats an error between two species in the same family as less severe than an error between species from different orders \cite{bertinetto2020making}.

Existing hierarchical classification methods can be broadly divided into single-level and full-taxonomy prediction approaches. In a bottom-up single-level approach, the model predicts the finest label, and higher-level labels are inferred from the taxonomy. This guarantees a valid taxonomy path, but it relies heavily on the correctness of the fine-grained prediction. In a top-down approach, the model first predicts a coarse label and then refines the prediction toward finer labels. This can be useful when coarse-level recognition is easier, but early mistakes may propagate downward \cite{park2025visually}.

Full-taxonomy prediction aims to predict all hierarchy levels simultaneously. This is the setting used in this thesis. It provides explicit supervision at each level and allows the model to evaluate performance at order, family, and species levels separately. However, because different heads may predict different levels independently, full-taxonomy prediction must address the problem of taxonomy consistency \cite{park2025visually}.

\subsection{Taxonomy Consistency and Multi-Granularity Prediction}

A central challenge in hierarchical classification is maintaining consistency among predictions at different levels. A prediction is taxonomically consistent if the predicted labels form a valid path in the hierarchy. For example, a predicted species should belong to the predicted family, and that family should belong to the predicted order. If the species and family predictions do not match the known taxonomy, the model produces an inconsistent prediction path.

This issue is especially important in multi-granularity fine-grained recognition. Different levels of the hierarchy may require different visual evidence. Coarse-level labels often depend on broad object-level cues, while fine-grained species labels may depend on subtle local details. H-CAST emphasizes this problem by showing that classifiers at different hierarchy levels may attend to different image regions, leading to inconsistent predictions even when semantic supervision is used \cite{park2025visually}.

Several metrics are used to evaluate hierarchical prediction. Level-wise accuracy measures the classification accuracy at each individual level, such as order accuracy, family accuracy, and species accuracy. However, level-wise accuracy alone does not fully capture whether the predictions form a valid taxonomy path. Tree-based InConsistency Error rate (TICE) measures the proportion of samples whose predicted path is inconsistent with the hierarchy. Full-Path Accuracy (FPA), on the other hand, counts a prediction as correct only when all hierarchy levels are predicted correctly. Therefore, FPA evaluates both correctness and hierarchical agreement, while TICE focuses on the validity of the predicted path \cite{park2025visually}.

Multi-granularity prediction is useful because it allows a model to provide predictions at different levels of specificity. In practical bird recognition, an image may not always contain enough detail for confident species-level prediction, but a coarse-level family or order prediction may still be meaningful. At the same time, predicting multiple levels introduces a trade-off: the model must learn features that are useful for coarse categories while preserving subtle visual cues needed for fine-grained species discrimination.

\subsection{Representative Hierarchical Classification Methods}

A number of methods have been proposed to exploit hierarchy information in image classification. Bertinetto et al. studied how class hierarchies can be used to make ``better mistakes'' by encouraging errors that are closer to the ground-truth class in the hierarchy \cite{bertinetto2020making}. Their work highlights that not all classification errors are equally severe and motivates the use of hierarchy-aware evaluation and training objectives.

For fine-grained visual classification, Chang et al. proposed FGN, which reformulates traditional FGVC as a multi-granularity recognition problem \cite{chang2021your}. Instead of predicting only the finest category, FGN considers a coarse-to-fine label hierarchy and uses level-specific classification heads to reduce negative transfer across granularity levels. The method also allows fine-grained features to contribute to coarse-level prediction, reflecting the observation that fine-level visual details can help recognize coarser labels.

Chen et al. proposed the Hierarchical Residual Network (HRN) for hierarchical multi-granularity classification \cite{Chen_2022_CVPR}. HRN models hierarchical feature interaction by adding granularity-specific features from parent levels to child-level features through residual connections. It also uses a combinatorial loss based on the tree hierarchy to exploit label relationships and transfer hierarchical knowledge during training \cite{Chen_2022_CVPR}. This approach focuses on transferring hierarchical knowledge between levels and improving feature interaction across granularities.

With the rise of Vision Transformers, Wang et al. introduced TransHP, a hierarchical prompting method for image classification \cite{wang2023transhp}. TransHP learns prompt tokens representing coarse ancestor classes and injects the predicted coarse-class prompt into intermediate transformer features. The motivation is that ancestor-level information can guide the model to focus on subtle differences among descendant classes. Unlike residual CNN-based approaches, TransHP uses the transformer token mechanism to condition feature extraction on hierarchy information.

Another related direction focuses directly on prediction consistency. Wang et al. proposed Consistency-Aware Feature Learning (CAFL) for hierarchical fine-grained visual classification \cite{wang2023consistency}. CAFL encodes prediction consistency constraints through forward deduction and backward induction over the label hierarchy, aiming to improve both classification accuracy and consistency. This type of method shows that enforcing consistency is important, but it mainly focuses on semantic relations between labels and prediction scores.

More recently, H-CAST introduced a visually consistent approach to hierarchical image classification \cite{park2025visually}. Unlike methods that rely mainly on label relations or semantic losses, H-CAST builds on CAST and uses superpixel-based hierarchical visual grouping. It applies fine-to-coarse supervision at different stages of the model so that finer visual segments support fine-grained labels and coarser segments support higher-level labels. In addition, H-CAST uses Tree-Path KL Divergence to encourage semantic consistency across hierarchy levels.

H-CAST is the most directly related baseline for this thesis. It provides a strong image-only framework for hierarchical bird classification and addresses both visual and semantic consistency. However, in its CUB setting, H-CAST uses the order-family-species taxonomy but does not incorporate the CUB visual attribute annotations as explicit auxiliary supervision. This thesis builds on the H-CAST framework by introducing attribute-guided learning at the species-level branch. The goal is to investigate whether semantic attribute supervision and hard-sibling-aware attribute selection can further improve fine-grained species discrimination while preserving hierarchical prediction.